% Wymaga: booktabs. Nie wymaga pakietu float.
\makeatletter
\@ifundefined{tablehere}{\newenvironment{tablehere}{\par\vskip\intextsep\def\@captype{table}\begin{center}}{\end{center}\par\vskip\intextsep}}{}
\makeatother


\subsection{Optymalizacja front-endu}
\subsubsection{Opis optymalizacji}

W pierwotnej wersji serwisu przeglądarka musiała pobrać i wykonać cztery biblioteki:
Bootstrap (CSS ok.\ 280\,KB, JS ok.\ 208\,KB), Font Awesome (CSS ok.\ 130\,KB
i czcionki), jQuery (285\,KB) i moment.js (183\,KB), z których strona używała
niewielkiego ułamka: kilku klas Bootstrapa, przełącznika menu, trzech ikon,
żądań JSON, wstawiania elementów i jednego formatowania daty. Wszystkie skrypty
w \texttt{<head>} blokowały parsowanie HTML, \texttt{analytics.js} dodatkowo
wykonywał aktywne oczekiwanie przez 100\,ms, a \texttt{style.css} dołączał
\texttt{theme.css} przez \texttt{@import} (dodatkowy, szeregowy odnośnik).
Animacja stopki była realizowana przez \texttt{setInterval} co 30\,ms,
który przez jQuery przepisywał \texttt{background} i \texttt{box-shadow}
(wymuszając przemalowanie, także w ukrytej karcie), a zdarzenie \texttt{mousemove}
nie było ograniczone. Treść wpisów była wstawiana jako ciąg HTML.

Zmiany:
\begin{itemize}
  \item usunięto wszystkie biblioteki (katalog \texttt{public/vendor});
    Bootstrap CSS zastąpiono własnym arkuszem (pasek nawigacji, kontener, siatka CSS Grid, karta),
    Bootstrap JS -- kilkoma liniami obsługi przełącznika menu, Font Awesome -- trzema ikonami
    jako inline SVG, jQuery -- \texttt{fetch} i DOM API, moment.js -- \texttt{Intl.DateTimeFormat},
  \item \texttt{theme.css} wbudowano w \texttt{style.css} (brak \texttt{@import}),
  \item skrypt \texttt{app.js} ładowany z \texttt{defer}, a \texttt{analytics.js} dopiero po zdarzeniu
    \texttt{load}, w czasie bezczynności przeglądarki (\texttt{requestIdleCallback});
    zdarzenia zgłoszone wcześniej trafiają do kolejki; usunięto aktywne oczekiwanie 100\,ms,
  \item dodano \texttt{<meta name="viewport">} i ikonę strony (unika zbędnego żądania \texttt{/favicon.ico} kończącego się 404),
  \item animacja stopki jako CSS \texttt{@keyframes} z \texttt{filter: hue-rotate()}
    na pseudo-elemencie, wstrzymywana poza ekranem (\texttt{IntersectionObserver}),
    w ukrytej karcie i przy \texttt{prefers-reduced-motion}; położenie kursora
    przekazywane zmienną CSS \texttt{--glow}, ograniczone do jednej aktualizacji na klatkę
    (\texttt{requestAnimationFrame}), z zapamiętanym \texttt{getBoundingClientRect()},
  \item treść wpisów wstawiana przez \texttt{textContent}, nie przez konkatenację HTML;
    pierwsze zdjęcie ładowane od razu, kolejne z \texttt{loading="lazy"}.
\end{itemize}

\subsubsection{Konfiguracja i parametry}

\begin{tablehere}
\centering
\begin{tabular}{l p{9cm}}
\toprule
\textbf{Było} & \textbf{Jest} \\
\midrule
\texttt{\$.getJSON(url, cb)} & \texttt{fetch(url).then(r => r.json())} \\
\texttt{\$('\#id').append(html)} & \texttt{createElement}, \texttt{textContent}, \texttt{appendChild} \\
\texttt{moment(d).format('D MMMM YYYY')} & \texttt{Intl.DateTimeFormat('pl-PL', \{day, month: 'long', year\})} \\
\texttt{data-bs-toggle="collapse"} & \texttt{classList.toggle('show')} \\
\texttt{<i class="fa-regular fa-calendar">} & inline \texttt{<svg>} z \texttt{fill="currentColor"} \\
\texttt{class="row g-4"}, \texttt{col-md-6} & \texttt{display:grid} z \texttt{@media (min-width: 768px)} \\
\texttt{setInterval(30 ms)} + \texttt{.css()} & \texttt{@keyframes} + \texttt{hue-rotate}, \texttt{IntersectionObserver} \\
\texttt{mousemove} + \texttt{text-shadow} z JS & zmienna CSS \texttt{--glow} ustawiana w \texttt{requestAnimationFrame} \\
skrypty w \texttt{<head>} bez atrybutów & \texttt{defer} (\texttt{app.js}), ładowanie po \texttt{load} (analityka) \\
\bottomrule
\end{tabular}
\caption{Zamienniki użytych funkcji}
\end{tablehere}

Zmiany widoczne dla użytkownika: daty są po polsku (np.\ „10 stycznia 2026” zamiast
„10 January 2026”), ikony są prostsze niż w Font Awesome, a pulsowanie cienia stopki
(zmienna wielkość) zastąpiono stałym cieniem.

\subsubsection{Wyniki pomiarów}

Rozmiary bez kompresji transportowej. Czcionki Font Awesome faktycznie pobierane
przez przeglądarkę to tylko \texttt{fa-regular-400.woff2}.

\begin{tablehere}
\centering
\begin{tabular}{l r r r}
\toprule
\textbf{Zasób} & \textbf{Przed [B]} & \textbf{Po [B]} & \textbf{Redukcja} \\
\midrule
CSS (Bootstrap, Font Awesome, \texttt{style.css}, \texttt{theme.css}) + czcionka & 430\,847 & 3\,695 & 99,1\% \\
JS (Bootstrap, jQuery, moment.js, \texttt{app.js}, \texttt{analytics.js}) & 679\,267 & 4\,149 & 99,4\% \\
\midrule
\textbf{Razem} & \textbf{1\,110\,114} & \textbf{7\,844} & \textbf{99,3\%} \\
\bottomrule
\end{tabular}
\caption{Rozmiary CSS, JS i czcionek przed i po optymalizacji}
\end{tablehere}

Liczba żądań o zasoby statyczne spadła z 10 (Bootstrap CSS, Font Awesome CSS, czcionka,
\texttt{style.css}, \texttt{theme.css}, trzy skrypty biblioteczne, \texttt{app.js},
\texttt{analytics.js}) do 3 (\texttt{style.css}, \texttt{app.js}, \texttt{analytics.js}),
przy czym tylko dwa pierwsze są potrzebne do pierwszego renderowania.

\begin{tablehere}
\centering
\begin{tabular}{l c c c c c}
\toprule
\textbf{Metryka} & \textbf{PSI mobile} & \textbf{PSI desktop} & \textbf{GTmetrix} & \textbf{Zmiana vs. stan wyjściowy} & \textbf{Uwagi} \\
\midrule
LCP [s]        & [—] & [—] & [—] & [—] & [—] \\
TBT [ms]       & [—] & [—] & [—] & [—] & [—] \\
CLS            & [—] & [—] & [—] & [—] & [—] \\
FCP [s]        & [—] & [—] & [—] & [—] & [—] \\
Speed Index [s]& [—] & [—] & [—] & [—] & [—] \\
\midrule
Wynik wydajności (0–100) & [—] & [—] & [—] & [—] & [—] \\
Rozmiar strony [MB]      & —   & —   & [—] & [—] & [—] \\
Liczba żądań HTTP        & —   & —   & [—] & [—] & [—] \\
Aktywność wątku głównego w bezczynności [\%] & — & — & [—] & [—] & [—] \\
\bottomrule
\end{tabular}
\caption{Wyniki pomiarów dla optymalizacji front-endu}
\end{tablehere}

\subsubsection{Wnioski}

Zasoby front-endowe skurczyły się z ok.\ 1,06\,MiB do ok.\ 7,7\,KiB. Przeglądarka nie
musi już pobierać, parsować i wykonywać czterech bibliotek przed pierwszym renderowaniem,
nie ma aktywnego oczekiwania ani synchronicznych skryptów w \texttt{<head>}, a w stanie
bezczynności animacja stopki nie obciąża wątku głównego (wyjątek: synchroniczny XHR w
\texttt{analytics.js}, który należy do optymalizacji sieci). Dodatkowo usunięto
podatność na wstrzyknięcie HTML przez treść wpisów.

Koszt to utrzymywanie własnego CSS i kodu pośredniczącego zamiast sprawdzonych bibliotek
oraz drobne różnice wizualne. Dla tak małej, jednostronicowej aplikacji jest to
opłacalne. Nie wykonano minifikacji ani pakowania (bundling), więc dalszy zysk jest możliwy
po dodaniu narzędzia budującego; własne pliki mają niewielki rozmiar.
